% BHU PYQ -- Midterm Examination 2025-2026
% Paper: MATMJ44 : Mechanics
% Exam: B.Sc. Year II Semester II Midterm Examination, 2025-2026

\documentclass[11pt,a4paper]{article}
\usepackage[a4paper,top=2.5cm,bottom=2.5cm,left=2.8cm,right=2.8cm,headheight=14pt]{geometry}
\usepackage{amsmath,amssymb}
\usepackage[T1]{fontenc}\usepackage[utf8]{inputenc}
\usepackage{lmodern,microtype,fancyhdr,enumitem,hyperref,lastpage,parskip}

\hypersetup{colorlinks=true,urlcolor=black,linkcolor=black,
  pdftitle={MATMJ44 Mechanics Midterm -- BHU B.Sc. Sem II 2025-26}}
\pagestyle{fancy}\fancyhf{}
\renewcommand{\headrulewidth}{0.4pt}\renewcommand{\footrulewidth}{0.4pt}
\fancyhead[L]{\small\textit{Banaras Hindu University}}
\fancyhead[R]{\small\textit{MATMJ44: Mechanics (Midterm 2025-26)}}
\fancyfoot[C]{\small Page \thepage\ of \pageref{LastPage}}

\newlist{parts}{enumerate}{2}
\setlist[parts,1]{label=(\alph*),leftmargin=*,itemsep=5pt,topsep=3pt}
\setlist[parts,2]{label=(\roman*),leftmargin=*,itemsep=3pt,topsep=2pt}
\newcommand{\pts}[1]{\hfill\textit{[#1]}}

\begin{document}
\begin{center}
  {\large\bfseries BANARAS HINDU UNIVERSITY}\\[2pt]
  {\large\bfseries Midterm Examination 2025-2026}\\[4pt]
  {\normalsize Program: B.Sc. \hfill Year: II \hfill Semester: II}\\[2pt]
  {\normalsize Subject: Mathematics \hfill Subject Code: MATMJ44 (Mechanics)}\\[4pt]
  \rule{0.8\linewidth}{0.5pt}\\[6pt]
  {\normalsize Date: 26/03/2026 \hfill Start Time: 9:00 AM}\\[2pt]
  {\normalsize Time Duration: One Hour \hfill Max. Marks: 20}
\end{center}
\rule{\linewidth}{0.4pt}
\smallskip
\noindent\textit{Instruction: 1) Answer any Four questions. The figures on the right hand margin indicate marks.\\
2) Write your section on top right of the front page of your answer script.\\
3) Give proper justification for your answers. Marks will not be awarded for guess work.}
\bigskip

\noindent\textbf{Question 1.}
\begin{parts}
  \item Define concurrent forces. Find analytical expression for resultant of two concurrent forces.\pts{3}
  \item The angle of inclination between two forces $P$ and $Q$ is $\theta$. If $P$ and $Q$ be interchanged in position, show that the resultant will be turned through an angle $\phi$, where $\tan\frac{\phi}{2} = \frac{P-Q}{P+Q}\tan\frac{\theta}{2}$, ($P > Q$).\pts{2}
\end{parts}
\medskip\rule{\linewidth}{0.2pt}

\medskip\noindent\textbf{Question 2.}
\begin{parts}
  \item State and prove the necessary and sufficient conditions for the equilibrium of a system of coplanar forces acting on a rigid body.\pts{5}
\end{parts}
\medskip\rule{\linewidth}{0.2pt}

\medskip\noindent\textbf{Question 3.}
\begin{parts}
  \item Define like parallel forces and obtain resultant of two like parallel forces. The resultant of two like parallel forces $P$ and $Q$ passes through a point $O$; when $P$ is increased by $R$ and $Q$ by $S$, the resultant still passes through $O$, and also when $Q, R$ replace $P, Q$, respectively; show that $S = R - \frac{(Q-R)^2}{P-Q}$.\pts{5}
\end{parts}
\medskip\rule{\linewidth}{0.2pt}

\medskip\noindent\textbf{Question 4.}
\begin{parts}
  \item Define moment of a force about a point. Show that the algebraic sum of moments of two parallel unlike forces about any point in their plane is equal to the moment of their resultant about that point.\pts{5}
\end{parts}
\medskip\rule{\linewidth}{0.2pt}

\medskip\noindent\textbf{Question 5.}
\begin{parts}
  \item Show that a system of particles at rest when acted upon by a system of coplanar forces remains in equilibrium if and only if the resultant work done by the forces is zero.\pts{5}
\end{parts}

\vfill\rule{\linewidth}{0.4pt}
\begin{center}\small
\href{https://bhu-exam.vercel.app/}{Click here to visit our website}\\[4pt]
\href{https://me-aryan.vercel.app/}{Click here to visit our contributor}\\[4pt]
\href{https://quashan-me.vercel.app/}{Click here to visit our contributor}
\end{center}
\end{document}
